% Experiment report. Fixed sections: see experiments/TEMPLATE/report.tex and experiments/README.md.
% Build: npm run experiment:report -- 32-rl-econ-research-additive
\documentclass{../polyexp}
\begin{document}
\input{generated}
\ExpHeader

\section{Question and hypothesis}
\textbf{The idea (David, 28 Nov 2026, issue \#166).} Experiments 29 and 31 each learned a schedule for one weight of
r15: economy down (to $\times0.25$--$\times0.5$ from turn 11, $+195$ against r15) and research up ($\times1.33$--$\times2$,
$\times2$ from turn 14, $+152$). Do the two gains add up when the learner plays both schedules at once? Both shift
stars from builds to techs (cutting economy to $\times0.5$ raised techs from 19.3 to 20.4 in 29's check), so they may
overlap or get in each other's way. The answer decides the next step: if they add up, learning one factor on top of
the other is enough; if not, the factors interact and need a joint state.

\textbf{Seen before the prediction} (stated so it is not mistaken for a blind prediction): the check that the new
pair encoding reproduces 29's and 31's games (below) also played the combined schedule on the same 300
validation seeds. It scored $-87$ against r15 there (148 better, 144 worse), where 29 and 31 alone scored $+175$ and
$+200$. At 300 seeds the standard error of a mean difference is about 100 points.

\textbf{Prediction} (after that look, before any run on the evaluation seeds):
\begin{enumerate}
  \item \emph{Not additive, and worse than either alone.} The combined schedule scores at least 100 points below
        29's economy policy and below 31's research policy; both comparisons go against it (no PASS).
  \item \emph{About r15's level.} Against r15 the combined schedule is between $-150$ and $+100$, no PASS.
\end{enumerate}
\textbf{Against the prediction:} a combined score above 29's alone would mean the gains add up at least in part,
and the 300-seed look was noise.

\section{Mechanism}
A new pair level set \texttt{econ-research25} (\texttt{rl\_value::strategic}): the 25 pairs of an economy level
(econ5: $\times0.25$, $\times0.5$, $\times0.75$, $\times1$, $\times1.5$) and a research level (research5: $\times0.5$,
$\times0.75$, $\times1$, $\times1.33$, $\times2$), index economy $\cdot 5$ + research. Both factors apply to r15's
weights; $(\times1, \times1)$ is r15 bit for bit, and each part moves one step at most per learner turn (tests check
both). \texttt{combine.py} writes 29's and 31's learned paths (state (turn, level), $m = 500$; with no income in the
state each policy is a fixed path of one level per turn) as level policies over the pair set:
\begin{itemize}
  \item \texttt{policy-combined.tsv}: 29's economy path and 31's research path together (the candidate);
  \item \texttt{policy-econ-only.tsv}, \texttt{policy-research-only.tsv}: one path, the other part at $\times1$. On
        300 validation seeds they play exactly 29's and 31's games, seed for seed (the encoding check).
\end{itemize}
No new training: this experiment tests a combination of two learned policies, not a new policy.

\section{Setup}
As experiments 29 and 31: learner base \texttt{greedy-full-probe-perf-r15}; opponents r15, probe-perf and the anchor
probe-perf, shuffled per seed; Perfection, 16$\times$16 Continents, 4 players, Crazy, learner Xin-xi in the human
seat. Evaluation seeds $10^9\dots$ (10{,}000), the same as 29 and 31.

Commands (from the repo root):
\begin{verbatim}
python3 experiments/32-rl-econ-research-additive/combine.py
E=experiments/32-rl-econ-research-additive
rl-gate --candidate rl-policy:$E/policy-combined.tsv --base greedy-full-probe-perf-r15
rl-gate --candidate rl-policy:$E/policy-combined.tsv --base rl-policy:experiments/29-rl-econ-level/policy-b0-m500.tsv --no-info15
rl-gate --candidate rl-policy:$E/policy-combined.tsv --base rl-policy:experiments/31-rl-research-level/policy-b0-m500.tsv --no-info15
\end{verbatim}

\section{Results}
\ExpResults

All numbers are copied from the \texttt{rl-gate} JSONs in this folder (28 Nov 2026; per-seed rows dropped, the runs are
deterministic). Encoding check (300 validation seeds $7\cdot10^8\dots$): \texttt{policy-econ-only.tsv} and
\texttt{policy-research-only.tsv} give exactly the scores of 29's and 31's policies on every seed.

\paragraph{Evaluation seeds} (10{,}000, \texttt{rl-gate}):

\begin{center}\small
\begin{tabular}{llrrrl}
\toprule
candidate & base & mean $d$ (99\% bounds) & better / worse & sign $p$ & verdict \\
\midrule
combined & r15         & $-27$ ($-74$, $+19$)    & 4{,}999 / 4{,}499 & $2\cdot10^{-7}$ & PASS (sign test) \\
combined & 29 economy  & $-222$ ($-261$, $-184$) & 3{,}936 / 5{,}059 & 1.0 & FAIL (mean worse) \\
combined & 31 research & $-180$ ($-221$, $-138$) & 3{,}966 / 5{,}192 & 1.0 & FAIL (mean worse) \\
\bottomrule
\end{tabular}
\end{center}

Score distribution on the same seeds (best of 10: the expected best of 10 random seeds):

\begin{center}\small
\begin{tabular}{lrrrrrr}
\toprule
learner & mean & p10 & median & p90 & best of 10 & best \\
\midrule
r15                  & 7{,}199 & 1{,}480 & 8{,}655 & 10{,}785 & 11{,}441 & 18{,}140 \\
29 economy           & 7{,}394 & 1{,}535 & 9{,}110 & 10{,}740 & 11{,}291 & 16{,}345 \\
31 research          & 7{,}351 & 1{,}520 & 9{,}130 & 10{,}755 & 11{,}307 & 16{,}785 \\
combined             & 7{,}172 & 1{,}550 & 8{,}955 & 10{,}590 & 11{,}085 & 15{,}445 \\
\bottomrule
\end{tabular}
\end{center}

15 opponents (information only), combined against r15: $-28$ ($-104$, $+47$), 1{,}587 better / 1{,}238 worse.

\section{Behaviour}
The combined learner plays both paths as written (the eval's factor-by-turn output matches \texttt{combine.py}'s
paths exactly): economy at $\times0.25$--$\times0.5$ from turn 11 and research at $\times2$ from turn 14, i.e.\ in the
second half of the game it builds a quarter to a half of r15's economy and researches at twice r15's weight.

It wins more seeds than r15 (53\% of the differing seeds) and has a better lower tail (p10 1{,}550), but its strong
games are weaker: p90 $-195$, best of 10 $-356$, best game $-2{,}695$. Each single schedule already trades a little of
the top for the middle (29 and 31: best of 10 about $-140$ against r15); the combination does it twice, and loses
more at the top than it gains in the middle.

\section{Interpretation}
\textbf{Prediction 1 (worse than either alone by at least 100, no PASS against them): right} ($-222$ and $-180$, both
FAIL on the mean guard).
\textbf{Prediction 2 (between $-150$ and $+100$ against r15, no PASS): half right.} The mean is $-27$, inside the range,
but the sign test passes: the combination is better than r15 on more seeds, and worse by more on the others.

\textbf{The gains do not add up; together they cancel.} 29's and 31's schedules each move the same thing, the share
of stars that goes to techs rather than builds, and each was learned with the other weight at r15's value. Played
together, the shift is about doubled, and it goes past the best split. The single-factor checks showed this shape
already: economy $\times0.5$ every turn gained $+309$ but $\times0.25$ only $+60$; research $\times1.33$ gained $+89$
but $\times2$ $-36$. Each factor has an optimum with a steep fall beyond it, and the combination lands beyond both.
Where it falls is in the strong games: those are the ones where a growing economy pays most, and the combined
learner gives up economy and pushes research at the same time.

\textbf{What this means for the method.} A schedule learned for one factor is a best response to the other factors
at their old values. It cannot be combined with another such schedule as if the factors were independent: the
factors interact through the shared budget. Learning them jointly, or one on top of the other, is required.

\section{Threats to validity}
\begin{itemize}
  \item \emph{The look before the prediction.} The 300-seed sanity run showed $-87$ before the predictions were
        written (stated in the question). The evaluation seeds were not seen.
  \item \emph{One combination.} Only the two paths as learned were combined; a weaker version (e.g.\ economy $\times0.5$
        with research $\times1.33$) might add up. This experiment shows that the learned paths do not, not that no
        combination does.
  \item \emph{Mechanism not measured directly.} The star split (cities, income, techs per game) is not in the gate
        output; the overshoot reading rests on the single-factor checks of 29 and 31 and the score distribution.
\end{itemize}

\section{Verdict and next step}
\textbf{Not kept.} The combination is about r15's level ($-27$; the sign test passes, the mean does not move) and
far below either schedule alone ($-222$ against 29, $-180$ against 31). The economy and research gains of 29 and 31
are not additive: the factors interact through the shared star budget.

Next step (one idea): learn the two levels jointly, state (turn, economy level, research level), with a random walk
over both (\texttt{econ-research25}, \texttt{--driver walk}). That needs \texttt{fit-level} to handle 25 levels (its
per-action arrays are sized for five). The cheaper alternative is to learn one factor on top of the other: research
levels with 29's economy schedule built into the base.

\section{Notes for the paper}
Draft claim: two single-factor schedules that each pass against r15 ($+195$ economy, $+152$ research) cancel when
played together (combined $-27$ against r15, $-222$ against economy alone), because both shift the same star budget
and each was learned with the other weight fixed. This motivates the joint state. Carried by the evaluation table
and the score-distribution table.
\todo{jakob}{Is this worth a place in the paper? Where?}

\end{document}
